\thispagestyle{empty}

\begin{center}
    \textbf{\large ABSTRACT}
\end{center}

\vspace{14pt}

\begin{singlespace}
\noindent\justifying
This paper presents the development of WeFIND, a collaborative mobile application designed to support pet owners and the community in disseminating information and searching for lost and found animals. The project was based on field research and a comparative analysis of similar platforms, which helped identify difficulties related to scattered information, locating cases, and the exposure of contact details. The findings informed the definition of requirements and UML models that guided development. The application brings together geolocation, the consultation and filtering of reports, sighting registration, private communication, and the generation of QR-code posters for sharing or printing. Development uses React Native and Expo for the mobile environment and the Supabase platform with a PostgreSQL relational database for authentication, storage, and data synchronization. Exploratory manual functional checks and a usability evaluation with end users indicated a positive perception of ease of use, considering the limitations of the sample and procedures.

\vspace{18pt}

\noindent
\textbf{Keywords:} Lost animals. Mobile applications. Geolocation. React Native. Supabase. Collaborative systems.
\end{singlespace}

\clearpage
